\documentclass[11pt]{article}
\usepackage[margin=1in]{geometry}
\usepackage{amsmath,amssymb,amsthm}
\usepackage{hyperref}
\usepackage{listings}
\usepackage{xcolor}
\usepackage{booktabs}
\usepackage{array}
\usepackage{microtype}
\usepackage{cite}
\usepackage{url}
\usepackage{skak}        % chess symbols
\usepackage{chessboard}  % board diagrams from FEN

% Allow line breaks at underscores and dots in \texttt
\renewcommand{\texttt}[1]{%
  \begingroup\ttfamily\hyphenchar\font=`\-\relax
  \def\_{\textunderscore\discretionary{}{}{}}%
  #1\endgroup}

% Global stretch to prevent overfull boxes from long identifiers
\emergencystretch=3em

\newtheorem{theorem}{Theorem}
\newtheorem{lemma}[theorem]{Lemma}
\newtheorem{proposition}[theorem]{Proposition}
\newtheorem{corollary}[theorem]{Corollary}
\theoremstyle{definition}
\newtheorem{definition}[theorem]{Definition}
\newtheorem{finding}[theorem]{Finding}
\theoremstyle{remark}
\newtheorem{remark}[theorem]{Remark}

\lstset{
  basicstyle=\ttfamily\small,
  keywordstyle=\color{blue},
  commentstyle=\color{gray},
  breaklines=true,
  frame=single,
  xleftmargin=1em,
  escapeinside={(*}{*)},
}

\title{Yanasse: Finding New Proofs from Deep Vision's Analogies\\[0.5em]
\large Part~2: Probability Theory $\to$ Set Theory}
\author{Alexandre Linhares\thanks{alexandre@linhares.ltd. Affiliations: Getulio Vargas Foundation and ARGO LABS}}
\date{April 2026}

\begin{document}
\maketitle

\begin{abstract}
This is Part~2 of the Yanasse series, applying cross-area tactic
transfer to the pair $\textsc{Probability} \to
\textsc{SetTheory}$.  The system identifies tactic schemas that are
heavily used in probability theory but rare or absent in set
theory (cardinal and ordinal arithmetic), matches source and target
proof states via GPU-accelerated relational analogy on a MacBook Air
(Apple MPS backend), and asks an AI reasoning agent to semantically
adapt the source pattern.

The method produces 6~Lean-verified new proofs from 4 successful
schema transfers out of 10 attempts (40\% schema success rate).  The
proofs target two theorems in cardinal and ordinal arithmetic:
\texttt{mul\_eq\_left\_iff} (when cardinal multiplication absorbs:
$a \cdot b = a$) and \texttt{omega0\_lt\_preOmega\_iff} (characterizing
$\omega < \operatorname{preOmega}(x) \Leftrightarrow \omega < x$).

The deepest transfer is a \emph{proof-architecture transfer}: the
probability-theory pattern ``case-split on a structural niceness
predicate'' (e.g., \texttt{IsSFiniteKernel\,$\kappa$}) maps to the
set-theory pattern ``case-split on whether a cardinal is infinite''
(\texttt{$\aleph_0 \leq a$}).  The two predicates share no
identifiers, types, or namespaces; the connection is purely semantic.
Cross-pair comparison with Part~1 (Representation Theory) confirms the
(head, modifier) decomposition finding: the same three schemas that
transferred to representation theory---\texttt{ext1},
\texttt{any\_goals}, \texttt{by\_cases}---also transfer to set theory,
while the same three---\texttt{lift}, \texttt{measurability},
\texttt{congrm}---fail in both areas.
\end{abstract}

\newpage
\tableofcontents
\newpage

%======================================================================
\section{Introduction}\label{sec:intro}
%======================================================================

Part~1 of this series~\cite{yanasse_part1} demonstrated that proof
strategy patterns can be systematically transferred across distant areas
of Mathlib, producing Lean-verified new proofs at a 40\% success rate on
the pair Probability $\to$ Representation Theory.  The key finding was
the \emph{(head, modifier) decomposition}: tactic heads are
domain-gated, but modifiers carry domain-general patterns that transfer.

This paper applies the same method to a second pair:
$\textsc{Probability} \to \textsc{SetTheory}$.  Set theory in
Mathlib comprises cardinal arithmetic
(\texttt{Mathlib.SetTheory.Cardinal}), ordinal arithmetic
(\texttt{Mathlib.SetTheory.Ordinal}), and ordinal notation
(\texttt{Mathlib.SetTheory.Ordinal.Notation}).  The mathematical
content---transfinite induction, cardinal absorption laws, ordinal
fixed-point theory---is structurally distant from probability's
measure-theoretic world.

The pair is ranked second overall in the Yanasse pair-potential
metric (Section~\ref{sec:method}), and the target area
(\texttt{Mathlib.SetTheory}) contains 2{,}076 extracted proof states,
roughly three times the size of the Representation Theory target (672).

\subsection{Overview of the Deep Vision framework}

The Deep Vision framework~\cite{linhares_deep_vision,
linhares2024deepvision}, growing out of the Fluid Analogies Research
Group tradition~\cite{hofstadter_fluid_1995, hofstadter_go_1999,
hofstadter_surfaces_2013}, represents any structured domain as a
\emph{relational network}: entities connected by typed, directed
relations.  An analogy optimization finds the entity permutation that
maximizes relational consistency between two networks.  The matching
code (\texttt{deep\_vision\_lib.py}) is entirely domain-independent:
it receives two relational structures and returns a similarity score
and entity correspondence, without knowing which domain it is
processing.

\paragraph{Chess analogies.}  In chess~\cite{linhares_active_2005,
linhares_understanding_2007, linhares_questioning_2010,
linhares_entanglement_2012, linhares_what_2013,
linhares_emergence_2014}, entities are pieces and relations are attack,
defense, blocking, confinement, and pinning.  Two pieces in entirely
different positions can play ``the same role'' if their relational
profiles match.  The matcher discovers such correspondences across
positions that differ radically in surface features---12-piece endgames
matched against 5-piece positions, opening traps matched against sparse
endgames~\cite{yanasse_part1}.

\paragraph{From chess to proofs.}  The same matching engine, with only
the relation extractor swapped, handles Lean~4 proof states: entities
become hypotheses, goals, and types; relations become rewrite,
fit/apply, head-match, structure, equality, reflexive, witness,
bidirectional, kernel/simp, decidable, and lemma-needed (14 relation
types).  No chess knowledge is used; no proof-specific heuristics are
added to the matcher.  The proof-state matcher has previously been
used to discover a formal proof of Wolstenholme's
theorem~\cite{linhares_wolstenholme_2026} in Lean~4
(see also~\cite{linhares_sketch_2026}).

\subsection{Results preview}

Applying the method to Probability $\to$ Set Theory produces:

\begin{itemize}
\item 6 Lean-verified new proofs from 4 successful schema transfers
  out of 10 attempts (40\% schema success rate, with one schema
  producing 3 independent proof variants).
\item A \emph{proof-architecture transfer}: the deepest result
  reorganizes a cardinal-arithmetic proof around a top-level
  infinite/finite dichotomy, mirroring probability's
  s-finite/non-s-finite split (Section~\ref{sec:proof3}).
\item Cross-pair comparison confirming the (head, modifier)
  decomposition from Part~1 (Section~\ref{sec:cross_pair}).
\item 6 negative results with mathematical diagnoses
  (Section~\ref{sec:failures}).
\end{itemize}

%======================================================================
\section{Methodology}\label{sec:method}
%======================================================================

The methodology is identical to Part~1~\cite{yanasse_part1}.  We
summarize the pipeline and note the parameters specific to this pair.

\subsection{Pipeline summary}

\begin{enumerate}
\item \textbf{Schema extraction.}  Parse the Mathlib proof-state
  corpus (217{,}133 entries).  Each tactic invocation becomes a schema
  tuple $(\textit{head}, \textit{arity}, \textit{has\_with},
  \textit{uses\_lemma})$.

\item \textbf{$z$-score ranking.}  For each (area, schema) pair,
  compute the $z$-score of the schema's frequency across 27 areas.
  A schema with $z \geq +2$ in the source and $z \leq -1$ (or absent)
  in the target is a transfer candidate.

\item \textbf{GPU-accelerated analogy matching.}  Sample 8 source
  proofs per schema, match each against all 2{,}076 Set Theory proof
  states using the \texttt{BatchMatcher} on MPS.  Retain the top-5
  target analogues by normalized score.

\item \textbf{Semantic adaptation.}  A fresh AI reasoning session
  (Claude Opus, 50-minute budget) adapts the source tactic's
  \emph{mathematical strategy}---not its syntax---to the target
  theorem.

\item \textbf{Lean verification.}  The adapted proof is compiled
  with \texttt{lake env lean}.  Exit code~0 with no \texttt{sorry}
  constitutes verification.
\end{enumerate}

\subsection{Pair-specific parameters}

\begin{itemize}
\item \textbf{Source:} \texttt{Mathlib.Probability} (30{,}384 tactic
  invocations, the largest Mathlib area).
\item \textbf{Target:} \texttt{Mathlib.SetTheory} (2{,}076 proof
  states---3$\times$ larger than the Representation Theory target).
\item \textbf{Matching cost:} $8 \times 2{,}076 = 16{,}608$
  individual analogy computations per schema, completing in
  approximately 15--60~minutes on MPS depending on source proof-state
  size.
\item \textbf{Top-10 transfer candidates} (by gap score):
  \texttt{filter\_upwards}, \texttt{congr\,with}, \texttt{fun\_prop},
  \texttt{ext1}, \texttt{lift}, \texttt{any\_goals},
  \texttt{measurability}, \texttt{by\_cases}, \texttt{congrm},
  \texttt{push}.
\end{itemize}

Note that the same 10 schemas appear as in Part~1 (with the addition
of \texttt{push}, replacing \texttt{all\_goals} which was already
tested against Representation Theory).  This overlap enables
direct cross-pair comparison.

%======================================================================
\section{Results: Probability $\to$ Set Theory}\label{sec:results}
%======================================================================

\begin{table}[h]
\centering\small
\caption{Transfer results for Probability $\to$ Set Theory.}
\label{tab:results}
\begin{tabular}{@{}clcccp{4.5cm}@{}}
\toprule
Item & Schema head & \texttt{with} & Ar.\ & Result & Adaptation \\
\midrule
10 & \texttt{filter\_upwards} & \checkmark & 1 & Failed & No \texttt{Filter} in ordinal arith.\ \\
11 & \texttt{congr} & \checkmark & 0 & Failed & No extensionality for Cardinal \\
12 & \texttt{fun\_prop} & & 0 & Failed & Domain-gated (Measurable) \\
13 & \texttt{ext1} & & 1 & \textbf{New proof} & \texttt{inductionOn + mk\_congr} \\
14 & \texttt{lift} & \checkmark & 5 & Failed & No \texttt{CanLift} for Ordinal \\
15 & \texttt{any\_goals} & & 1 & \textbf{New proof} & \texttt{any\_goals lia} \\
16 & \texttt{measurability} & & 0 & Failed & Domain-gated ($\sigma$-alg.) \\
17 & \texttt{by\_cases} & & 4 & \textbf{New proof} & \texttt{by\_cases ha : $\aleph_0 \leq$ a} \\
18 & \texttt{congrm} & & 5 & Failed & No shared outer operators \\
19 & \texttt{push} & & 3 & \textbf{New proofs\,(3)} & \texttt{rwa [$\leftarrow$ preOmega\_omega0]} \\
\bottomrule
\end{tabular}
\end{table}

\noindent Four transfer attempts succeeded, producing 6 Lean-verified
new proofs (item~19 yielded 3 independent proof variants).  The
target theorems are:

\begin{enumerate}
\item[\textbf{P1.}] \textbf{Cardinal multiplication absorption}
  (\texttt{mul\_eq\_left\_iff}, items~13, 15, 17).
  For cardinals $a, b$:
  \[
    a \cdot b = a \;\Longleftrightarrow\;
    \bigl(\max(\aleph_0, b) \leq a \;\wedge\; b \neq 0\bigr)
    \;\vee\; b = 1 \;\vee\; a = 0.
  \]
  Three new proofs via three different transferred schemas.

\item[\textbf{P2.}] \textbf{Pre-omega comparison}
  (\texttt{omega0\_lt\_preOmega\_iff}, item~19).
  For ordinals:
  \[
    \omega < \operatorname{preOmega}(x)
      \;\Longleftrightarrow\; \omega < x.
  \]
  Three proof variants via the transferred ``push'' normalization
  pattern.
\end{enumerate}

\subsection{New proof P1a: \texttt{ext1} $\to$
  \texttt{Cardinal.inductionOn + mk\_congr}}\label{sec:proof1a}

\paragraph{Theorem.}
For cardinals $a$ and $b$, the absorption law $a \cdot b = a$
holds if and only if $\max(\aleph_0, b) \leq a$ and $b \neq 0$,
or $b = 1$, or $a = 0$.

\paragraph{Context.}
Cardinal arithmetic diverges sharply from natural-number arithmetic
once infinite cardinals appear.  The absorption law
$a \cdot b = a$ when $a$ is infinite and $b \leq a$ was first
established (for $a \cdot a = a$) by Hessenberg in
1906~\cite{hessenberg_1906}, later generalized by Tarski's theorem
on cardinal products~\cite{tarski_1925} and systematized in
Jech~\cite{jech_set_theory_2003} (Chapter~5) and
Kunen~\cite{kunen_set_theory_2011} (Chapter~I, \S10).
Holz, Steffens, and Weitz~\cite{holz_steffens_weitz_1999} provide
a modern treatment in the context of Shelah's pcf
theory~\cite{shelah_cardinal_1994}, which reveals deeper structure
behind singular cardinal arithmetic.  The characterization
in \texttt{mul\_eq\_left\_iff}---providing the complete equivalence
including the finite and degenerate cases---appears in
\texttt{Mathlib.SetTheory.Cardinal.Arithmetic}, whose formalization
follows Jech's treatment.  The result is a workhorse lemma for
cardinal exponentiation and Galvin--Hajnal-type
bounds~\cite{galvin_hajnal_1975}.

\paragraph{New proof (ext1 analog).}
The original proof dispatches the two ``easy'' backward cases
($b = 1$ and $a = 0$) with \texttt{all\_goals simp}, which invokes
the simplification lemmas \texttt{mul\_one} and \texttt{zero\_mul}.
The transferred proof replaces this automation with an
\emph{explicit type-equivalence witness}, the set-theoretic analog
of the extensionality principle:

\begin{itemize}
\item For $a \cdot 1 = a$: lift $a$ to a concrete type $\alpha$ via
  \texttt{Cardinal.inductionOn}; express $\#\alpha \cdot 1$ as
  $\#(\alpha \times \mathtt{PUnit})$ via \texttt{Cardinal.mul\_def};
  close with \texttt{Cardinal.mk\_congr (Equiv.prodPUnit\,$\alpha$)},
  the canonical equivalence $\alpha \times \mathtt{PUnit} \simeq \alpha$.

\item For $0 \cdot b = 0$: analogously, lift $b$ and use
  \texttt{Equiv.pemptyProd\,$\beta$} for
  $\mathtt{PEmpty} \times \beta \simeq \mathtt{PEmpty}$.
\end{itemize}

\paragraph{Source pattern.}  In probability theory, \texttt{ext1\,$\omega$}
reduces set equality $S = T$ to pointwise membership
$\forall \omega,\; \omega \in S \leftrightarrow \omega \in T$---a
workhorse for measurability proofs~\cite{yanasse_part1}.

\paragraph{Adaptation.}  Direct \texttt{ext1} fails: Cardinals have no
\texttt{@[ext]} lemma (they are quotient types, not structures with
component-wise decomposition).  The semantic analog of ``reduce
equality to constituent-level witness'' is: lift to a concrete
representative type via \texttt{inductionOn}, then exhibit an
explicit type equivalence via \texttt{mk\_congr}.

\paragraph{Analogy matching.}  The source proof
(\texttt{measurable\_stoppedValue}) had 51 entities; the top-3
target analogues were all tied at QAP score 18.19, matched from
\texttt{Cardinal/Arithmetic.lean}, \texttt{Cardinal/Basic.lean}, and
\texttt{Ordinal/FixedPoint.lean}.

\subsection{New proof P1b: \texttt{any\_goals} $\to$
  \texttt{any\_goals lia}}\label{sec:proof1b}

\paragraph{Theorem.}
Same as P1a: \texttt{mul\_eq\_left\_iff}.

\paragraph{New proof.}
The transferred tactic \texttt{any\_goals lia} replaces the
original \texttt{all\_goals simp} at the cleanup step of the
backward direction.  After the case-split
\texttt{rintro (...\,|\,rfl\,|\,rfl)}, two sub-goals remain:
$a \cdot 1 = a$ (from $b = 1$) and $0 \cdot b = 0$ (from $a = 0$).
The linear integer arithmetic decision procedure \texttt{lia}
closes both via internal normalization on Cardinals.

\paragraph{Source pattern.}  In
\texttt{Mathlib.Probability.Kernel.IonescuTulcea.Traj}, the tactic
\texttt{any\_goals lia} dispatches trivial arithmetic side-conditions
after a case split on trajectory indices.  It is a \emph{partial-sweep
combinator}: broadcast a decision procedure across all sub-goals,
close the easy ones, leave the hard ones for subsequent tactics.

\paragraph{Adaptation.}  The transfer is \emph{literal}: the source
tactic \texttt{any\_goals lia} compiles directly in the target
context, closing the same kind of arithmetic identities
($a \cdot 1 = a$, $0 \cdot b = 0$) that its probability-theory
counterpart closes (index arithmetic like $i \leq k + 1$).

\paragraph{Contrast with Part~1.}  In Part~1,
\texttt{any\_goals rfl} transferred to Representation Theory for
the \emph{same architectural reason}: after a case-split, one branch
is definitionally trivial.  The inner tactic differs (\texttt{rfl}
vs.\ \texttt{lia}), but the proof pattern---``broadcast a simple
closer across heterogeneous sub-goals''---is identical.

\subsection{New proof P1c: \texttt{by\_cases} $\to$ top-level
  infinite/finite dichotomy}\label{sec:proof3}

\paragraph{Theorem.}
Same as P1a: \texttt{mul\_eq\_left\_iff}.

\paragraph{Context.}
This is the deepest semantic transfer in the paper.  The original
Mathlib proof of \texttt{mul\_eq\_left\_iff} splits the biconditional
first (\texttt{refine $\langle$fun h $\Rightarrow$ ?, ?$\rangle$}),
then uses \texttt{rcases le\_or\_gt $\aleph_0$ a} deep inside the
forward direction.  The transferred proof reorganizes the entire
argument around a top-level \texttt{by\_cases} on whether $a$ is
infinite.

\paragraph{New proof.}
\begin{enumerate}
\item \texttt{by\_cases ha : $\aleph_0 \leq$ a} --- the structural
  ``niceness split.''
\item \textbf{Infinite case} ($\aleph_0 \leq a$): Cardinal
  multiplication absorbs: $a \cdot b = \max(a, b)$ by
  \texttt{mul\_eq\_max\_of\_aleph0\_le\_left}.  From
  $a \cdot b = a$ derive $b \leq a$ via the max characterization.
  The backward direction dispatches the three disjuncts with
  \texttt{mul\_eq\_left}, \texttt{simp}, \texttt{simp}.
\item \textbf{Finite case} ($a < \aleph_0$): The first disjunct
  (requiring $\aleph_0 \leq a$) is impossible.  A nested
  \texttt{by\_cases h2a : a = 0} handles the sub-degenerate case.
  When $a \neq 0$, show $b$ must also be finite (otherwise
  $a \cdot b \geq b \geq \aleph_0 > a$), cast to $\mathbb{N}$
  via \texttt{lt\_aleph0}, and use \texttt{Nat.mul\_eq\_left}.
\end{enumerate}

\paragraph{Source pattern.}  In
\texttt{Mathlib.Probability.Kernel.Composition.MeasureComp}, the
theorem \texttt{comp\_compProd\_comm} uses:
\begin{lstlisting}
by_cases h(*$\kappa$*) : IsSFiniteKernel (*$\kappa$*)
\end{lstlisting}
This is a \emph{structural niceness guard}: s-finite kernels support
composition-product theory; non-s-finite kernels make
\texttt{compProd} collapse to zero, so \texttt{simp} with
collapsing lemmas closes the negative case in one line.

\paragraph{Adaptation.}  The semantic mapping is:

\medskip
\begin{center}\small
\begin{tabular}{@{}lll@{}}
\toprule
& Source (Probability) & Target (Set Theory) \\
\midrule
Niceness predicate & \texttt{IsSFiniteKernel\,$\kappa$}
  & $\aleph_0 \leq a$ \\
What it enables & Composition-product theory
  & Max-absorption arithmetic \\
Negative case & \texttt{compProd} collapses $\to$ \texttt{simp}
  & Falls to $\mathbb{N}$ arithmetic \\
Positive case & Extensionality + rewrite
  & \texttt{mul\_eq\_max} + \texttt{max\_eq\_left} \\
\bottomrule
\end{tabular}
\end{center}
\medskip

\noindent The predicates share no identifiers, types, or namespaces.
The connection is semantic: both determine whether the main algebraic
theory of their respective domains applies.  This is a transfer at the
level of \emph{proof architecture}, not tactic syntax.

\paragraph{Cross-pair comparison.}  In Part~1, \texttt{by\_cases}
transferred to Representation Theory as ``case-split on $f = 0$''
(a degenerate-morphism guard).  The Set Theory adaptation is
arguably deeper: the probability-theory analog
(\texttt{IsSFiniteKernel}) and the set-theory analog
($\aleph_0 \leq a$) both partition their domain's objects into
a ``nice regime'' where the main algebraic theory works and a
``degenerate regime'' requiring qualitatively different reasoning.
The morphism-triviality split in Part~1 was mathematically
unnecessary (both branches closed identically); the
infinite/finite split here is essential.

\subsection{New proof P2: \texttt{push} $\to$ ordinal normalization
  through an embedding}\label{sec:proof2}

\paragraph{Theorem.}
For ordinals $x$:
\[
  \omega < \operatorname{preOmega}(x)
    \;\Longleftrightarrow\; \omega < x.
\]
Here $\operatorname{preOmega} : \mathrm{Ord} \hookrightarrow
\mathrm{Ord}$ is the order embedding enumerating initial ordinals,
and $\omega$ is the first infinite ordinal.

\paragraph{Context.}
The function $\operatorname{preOmega}$ enumerates the \emph{initial
ordinals}---the smallest ordinals of each cardinality---a
construction going back to Hartogs~\cite{hartogs_1915} and
systematized by von Neumann~\cite{von_neumann_1923}.  In modern
notation, $\operatorname{preOmega}(\alpha) = \omega_\alpha$ for
infinite $\alpha$, with the convention
$\operatorname{preOmega}(n) = n$ for finite $n$.  The fixed-point
property $\operatorname{preOmega}(\omega) = \omega$ reflects the
fact that $\omega$ is the $\omega$-th initial ordinal---the
smallest aleph fixed point, studied extensively in
Jech~\cite{jech_set_theory_2003} (Chapter~2) and
Kunen~\cite{kunen_set_theory_2011} (Chapter~I, \S7).
Drake~\cite{drake_1974} gives a thorough treatment of the aleph
hierarchy and its fixed-point structure.  The biconditional
\texttt{omega0\_lt\_preOmega\_iff} exploits the fixed-point property
and the strict monotonicity of $\operatorname{preOmega}$ as an
order embedding; it appears in
\texttt{Mathlib.SetTheory.Cardinal.Aleph}.  The original Mathlib
proof uses a targeted rewrite:
\texttt{conv\_lhs => rw [$\leftarrow$ preOmega\_omega0,
preOmega\_lt\_preOmega]}.

\paragraph{New proofs (3 variants).}
The transferred ``push'' normalization pattern---normalize by pushing a
structural wrapper through an operator---yields three proof variants:

\begin{enumerate}
\item \textbf{Constructor + rwa.}  Split the iff; in each direction,
  \texttt{rwa [$\leftarrow$ preOmega\_omega0,
  preOmega\_lt\_preOmega]} introduces
  $\operatorname{preOmega}(\omega)$ and cancels the paired embeddings.
  Each rewrite is applied exactly once, avoiding the simp loop
  (see below).

\item \textbf{Targeted rewrite.}  \texttt{nth\_rw 1 [show $\omega =
  \operatorname{preOmega}(\omega)$ from preOmega\_omega0.symm]; exact
  preOmega\_lt\_preOmega}.  Replace only the first $\omega$, then
  close with the embedding's monotonicity lemma.

\item \textbf{Calc chain.}  A two-step calculation making the
  ``push-through'' steps explicit:
  $\omega < \operatorname{preOmega}(x)
  \Leftrightarrow \operatorname{preOmega}(\omega) <
  \operatorname{preOmega}(x)
  \Leftrightarrow \omega < x$.
\end{enumerate}

\paragraph{Source pattern.}  In
\texttt{Mathlib.Probability.Process.Stopping}, the tactic
\texttt{push \_ $\in$ \_} normalizes set membership by pushing
$\in$ through set-builder notation, preimages, and indexed unions:
$\omega \in \{x \mid P\,x\} \leadsto P\,\omega$.

\paragraph{Adaptation.}  The source pushes $\in$ through syntactic
set-construction wrappers; the target pushes $<$ through
$\operatorname{preOmega}$, an order-embedding wrapper.  The key
lemmas play parallel roles:

\medskip
\begin{center}\small
\begin{tabular}{@{}ll@{}}
\toprule
Source (Probability) & Target (Set Theory) \\
\midrule
\texttt{Set.mem\_setOf\_eq} & \texttt{preOmega\_lt\_preOmega} \\
(unfold set-builder) & (cancel paired embeddings) \\
\texttt{Set.mem\_preimage} & \texttt{preOmega\_omega0} \\
(push through preimage) & (introduce embedding at fixed point) \\
\bottomrule
\end{tabular}
\end{center}
\medskip

\paragraph{A failed approach.}  The most natural ``push'' analog---
\texttt{simp only [$\leftarrow$ preOmega\_omega0,
preOmega\_lt\_preOmega]}---loops: the backward rewrite
$\omega \to \operatorname{preOmega}(\omega)$ and the forward
simplification that cancels paired embeddings create an infinite
cycle.  The source's \texttt{push \_ $\in$ \_} never loops because
set-builder unfolding is one-directional.  The \texttt{rwa}-based
adaptation avoids this by applying each rewrite exactly once.

%======================================================================
\section{Cross-Pair Comparison with Part~1}\label{sec:cross_pair}
%======================================================================

Seven of the ten schemas tested against Set Theory were also tested
against Representation Theory in Part~1.  This enables direct
comparison:

\begin{table}[h]
\centering\small
\caption{Schema transfer results across two target areas.}
\label{tab:cross_pair}
\begin{tabular}{@{}lccc@{}}
\toprule
Schema head & RepTheory (Part~1) & SetTheory (Part~2)
  & Category \\
\midrule
\texttt{filter\_upwards} & \textbf{New proof} & Failed
  & Domain-gated \\
\texttt{congr\,with} & \textbf{New proof} & Failed
  & Domain-general \\
\texttt{fun\_prop} & Failed & Failed
  & Domain-gated \\
\texttt{ext1} & Failed & \textbf{New proof}
  & Domain-general \\
\texttt{lift} & Failed & Failed
  & Domain-gated \\
\texttt{any\_goals} & \textbf{New proof} & \textbf{New proof}
  & Domain-general \\
\texttt{measurability} & Failed & Failed
  & Domain-gated \\
\texttt{by\_cases} & \textbf{New proofs\,(2)} & \textbf{New proof}
  & Domain-general \\
\texttt{congrm} & Failed & Failed
  & Domain-gated \\
\texttt{push} & --- & \textbf{New proofs\,(3)}
  & Domain-general \\
\bottomrule
\end{tabular}
\end{table}

\begin{finding}[Consistent transfer/failure across targets]\label{find:consistency}
Of the 7~schemas tested in both areas:
\begin{itemize}
\item \textbf{2 transfer in both:} \texttt{any\_goals},
  \texttt{by\_cases}.  These are domain-general combinators whose
  proof-architectural patterns (partial sweep, niceness guard) exist
  in every mathematical area.
\item \textbf{3 fail in both:} \texttt{fun\_prop},
  \texttt{lift}, \texttt{measurability}.  These are domain-gated
  tactics requiring measure-theoretic or topological type structures
  absent from both targets.
\item \textbf{2 differ:} \texttt{filter\_upwards} (succeeds in
  RepTheory via modifier transfer, fails in SetTheory where
  ordinal arithmetic has no Filter API),
  \texttt{ext1} (fails in RepTheory where categorical
  extensionality is structurally incompatible, succeeds in
  SetTheory where \texttt{Cardinal.inductionOn} provides a semantic
  analog).
\end{itemize}
\end{finding}

The consistency of the ``never-transfer'' group---\texttt{fun\_prop},
\texttt{lift}, \texttt{measurability}---across both targets
strengthens the evidence that these schemas are \emph{intrinsically}
domain-gated, not merely unlucky in their target matching.  Future
runs could skip them entirely for non-measure-theoretic targets,
saving ${\sim}30$\% of compute.

The two schemas that differ between targets (\texttt{filter\_upwards}
and \texttt{ext1}) illustrate that the (head, modifier) decomposition
is target-sensitive: the modifier of \texttt{filter\_upwards}
(``\texttt{[LIST] with $\omega$}'') finds a structural analog in
representation theory's extensionality-based reasoning but not in
ordinal arithmetic's algebraic reasoning, while \texttt{ext1}'s
principle (``reduce equality to constituent-level witness'') finds
a semantic analog in cardinal arithmetic's type-equivalence
infrastructure but not in categorical extensionality.

%======================================================================
\section{The (Head, Modifier) Decomposition: Updated Evidence}%
\label{sec:finding}
%======================================================================

Part~1 introduced the (head, modifier) decomposition.  The Set Theory
results provide additional evidence and one refinement.

\begin{finding}[Proof-architecture transfer]\label{find:architecture}
The deepest transfers are not at the tactic level but at the
\emph{proof-architecture} level: the decision to organize a proof
around a top-level case-split on a structural niceness predicate
(\texttt{by\_cases}, item~17) transfers across areas because the
\emph{architecture}---split on niceness, handle degenerate case
trivially, proceed with main theory in nice case---is domain-general,
even when the specific niceness predicates are domain-specific.
\end{finding}

The \texttt{by\_cases} transfer to Set Theory is qualitatively
different from its Part~1 counterpart.  In Part~1, the case-split
($f = 0$) was mathematically unnecessary (both branches closed
identically); the transfer demonstrated that the \emph{pattern}
works but not that it \emph{improves} the proof.  In Part~2, the
infinite/finite split is mathematically essential: the infinite
case uses max-absorption ($a \cdot b = \max(a, b)$) while the
finite case casts to $\mathbb{N}$ and uses \texttt{Nat.mul\_eq\_left}.
These are fundamentally different proof strategies, and the
transferred architecture makes this dichotomy explicit at the top
level rather than burying it inside the forward direction.

%======================================================================
\section{Failures and Their Diagnoses}\label{sec:failures}
%======================================================================

\begin{description}
\item[Item 10: \texttt{filter\_upwards}.]
  \texttt{filter\_upwards} requires \texttt{Filter.Eventually} goals.
  Set Theory's ordinal notation proofs (the top-1 target:
  \texttt{nf\_repr\_split'}) have goals of the form $\mathtt{NF}\,o
  \wedge \mathtt{repr}\,o = \omega \cdot \mathtt{repr}\,o' + m$---propositional
  conjunctions about ordinal notation, not filter membership.
  The 55-to-3 entity ratio (source to target) also produced a
  spurious QAP score.

\item[Item 11: \texttt{congr\,with}.]
  \texttt{congr\,with} requires an extensionality theorem for the
  goal type.  \texttt{Cardinal} has no \texttt{@[ext]} lemma
  (confirmed by Lean: ``No applicable extensionality theorem found
  for type Cardinal'').  Positive controls confirmed that
  \texttt{congr\,with\,n} \emph{does} work on ordinal function
  equality (\texttt{nfp\_add\_zero}) and \texttt{Cardinal.sum}---a
  genuine transfer exists at rank~3, but the QAP matcher could not
  discriminate it from the degenerate rank-1 mapping (37-to-3 entity
  ratio).

\item[Item 12: \texttt{fun\_prop}.]
  \texttt{fun\_prop} handles \texttt{Measurable}, \texttt{Continuous},
  and \texttt{Differentiable} goals.  Set Theory has zero such
  goal types.  Same root cause as Part~1 item~2.

\item[Item 14: \texttt{lift}.]
  \texttt{lift} requires \texttt{CanLift} instances (e.g.,
  $\mathbb{R}_{\geq 0}^\infty \to \mathbb{R}_{\geq 0}$).  No
  \texttt{CanLift} instance exists for \texttt{Ordinal} in Mathlib:
  ordinals have no canonical simpler base type.  Five Lean probes
  all failed at typeclass synthesis.  Same root cause as Part~1
  item~5.

\item[Item 16: \texttt{measurability}.]
  A $\sigma$-algebra predicate closure tactic.  The target theorem
  (\texttt{repr\_opow}) proves that ordinal notation representation
  preserves exponentiation---a purely algebraic statement with no
  measurability predicates.  Same root cause as Part~1 item~7.

\item[Item 18: \texttt{congrm}.]
  \texttt{congrm} requires both sides of an equality to share a
  common outer-operator skeleton (integrals, products, set
  membership).  The target goal
  $\omega < \operatorname{preOmega}(x) \Leftrightarrow \omega < x$
  is an atomic order comparison---no composite structure to
  decompose.  Five QAP-matched targets were identically scored
  (23.53), indicating no discriminative structural match.
  Same root cause as Part~1 item~9.
\end{description}

%======================================================================
\section{Discussion}\label{sec:discussion}
%======================================================================

\subsection{The higher proof yield}

Set Theory's schema success rate (4/10 = 40\%) matches
Representation Theory's, but the \emph{proof yield} is higher:
6~Lean-verified proofs versus~4.  This may reflect the target
area's richer library of arithmetic lemmas.
Three of the four successful transfers target
\texttt{mul\_eq\_left\_iff}, a theorem with a complex case structure
and multiple algebraic regimes.  Such theorems offer more ``surface
area'' for diverse proof strategies.  In contrast, Representation
Theory's successful transfers targeted two different theorems, each
with a simpler proof structure.

\subsection{Convergence to a single target}

Three of the four successful schemas (items~13, 15, 17) converge on
the same target theorem (\texttt{mul\_eq\_left\_iff}).  This is
partly an artifact of the QAP matcher---the 2{,}076-entry target pool
has a small number of high-scoring ``attractor'' entries---and partly
reflects that \texttt{mul\_eq\_left\_iff} is genuinely a rich theorem
with a case-analytic structure amenable to multiple proof strategies.

\subsection{Depth of transfer}

The transfers range from shallow (item~15: literal
\texttt{any\_goals lia} compiles directly) to deep (item~17:
the ``niceness guard'' architecture requires semantic reasoning about
the role of $\aleph_0 \leq a$ in cardinal arithmetic).  Item~19
is intermediate: the ``push normalization'' pattern is genuinely
structural, but the target theorem is already a simple rewrite.

\subsection{QAP entity-ratio degeneration}

Items~10 and~11 both suffered from degenerate QAP scores caused by
extreme entity-ratio mismatches (55-to-3, 37-to-3).  When the
source proof state has $10\times$ more entities than the target, the
matcher can achieve high scores by mapping a small cluster of source
entities to the target while ignoring the rest.  Item~11's positive
control (a genuine transfer at rank~3) was masked by the degenerate
rank-1 score.

This motivates a planned improvement: penalizing matches where the
entity ratio exceeds a threshold (e.g., $5:1$), so that the matcher
favors structurally comparable proof states.

\subsection{Cost}

The full pipeline for this pair consumed:
\begin{itemize}
\item ${\sim}$25 minutes of GPU time per schema (MPS) for analogy
  matching ($8 \times 2{,}076$ comparisons).
\item ${\sim}$8 hours of AI reasoning across 10 semantic adaptation
  sessions.
\item ${\sim}$\$100--150 in API cost.
\end{itemize}

%======================================================================
\section{Conclusion}\label{sec:conclusion}
%======================================================================

The Probability $\to$ Set Theory pair achieves a 40\% schema success
rate (4/10), producing 6 new Lean-verified proofs.  The key result is a
\emph{proof-architecture transfer}: the probability-theory pattern
``case-split on a structural niceness predicate'' maps to cardinal
arithmetic's infinite/finite dichotomy, demonstrating that
cross-area transfer can operate at the level of proof organization,
not just tactic invocation.

Cross-pair comparison with Part~1 confirms that domain-gated
schemas (\texttt{fun\_prop}, \texttt{lift}, \texttt{measurability})
consistently fail across targets, while domain-general
combinators (\texttt{any\_goals}, \texttt{by\_cases}) consistently
transfer.  This consistency enables predictive filtering: future
runs can skip known-doomed schemas for non-measure-theoretic targets.

Part~3 will apply the method to \texttt{Probability} $\to$
\texttt{Condensed}, the hardest target tested so far (condensed
mathematics eliminates the Filter and $\sigma$-algebra vocabulary that
probability-theory schemas rely on).

\section*{Acknowledgments}

Named after Horacio Hideki Yanasse.  The Deep Vision matching engine
was originally developed for chess cognition research; its application
to formal mathematics was enabled by the domain-independent design of
\texttt{deep\_vision\_lib.py}.  All computations were performed on a
MacBook Air (Apple M-series, 32\,GB unified memory) using PyTorch~2.8
on the MPS backend.

\begin{thebibliography}{99}

\bibitem{yanasse_part1}
A.~Linhares,
``Yanasse: Finding new proofs from Deep Vision's analogies --- Part~1:
Probability Theory $\to$ Representation Theory,''
ARGO LABS Report for Distribution, April 2026.

\bibitem{linhares2024deepvision}
A.~Linhares,
``Deep Vision,''
ARGO LABS Internal Report, 2026.

\bibitem{hofstadter_fluid_1995}
D.~R.~Hofstadter and the Fluid Analogies Research Group, \emph{Fluid Concepts \& Creative Analogies: Computer Models of the Fundamental Mechanisms of Thought}. New York: Basic Books, 1995.

\bibitem{hofstadter_go_1999}
D.~R.~Hofstadter, \emph{G\"odel, Escher, Bach: An Eternal Golden Braid}, 20th anniversary ed. New York: Basic Books, 1999.

\bibitem{hofstadter_surfaces_2013}
D.~R.~Hofstadter and E.~Sander, \emph{Surfaces and Essences: Analogy as the Fuel and Fire of Thinking}. New York: Basic Books, 2013.

\bibitem{linhares_active_2005}
A.~Linhares, ``An active symbols theory of chess intuition,'' \emph{Minds and Machines}, vol.~15, pp.~131--151, 2005.

\bibitem{linhares_understanding_2007}
A.~Linhares and P.~Brum, ``Understanding our understanding of strategic scenarios: What role do chunks play?'' \emph{Cognitive Science}, vol.~31, no.~6, pp.~989--1007, 2007.

\bibitem{linhares_questioning_2010}
A.~Linhares and A.~E.~T.~A.~Freitas, ``Questioning Chase and Simon's (1973) `Perception in Chess','' \emph{New Ideas in Psychology}, vol.~28, no.~1, pp.~64--78, 2010.

\bibitem{linhares_entanglement_2012}
A.~Linhares, A.~E.~T.~A.~Freitas, A.~Mendes, and J.~S.~Silva, ``Entanglement of perception and reasoning in the combinatorial game of chess,'' \emph{Cognitive Systems Research}, vol.~13, no.~1, pp.~72--86, 2012.

\bibitem{linhares_what_2013}
A.~Linhares and D.~M.~Chada, ``What is the nature of the mind's pattern-recognition process?'' \emph{New Ideas in Psychology}, vol.~31, no.~2, pp.~108--121, 2013.

\bibitem{linhares_emergence_2014}
A.~Linhares, ``The emergence of choice: Decision-making and strategic thinking through analogies,'' \emph{Information Sciences}, vol.~259, pp.~36--56, 2014.

\bibitem{linhares_deep_vision}
A.~Linhares, ``Deep Vision: Seeing the essence of proofs through analogy,'' ARGO LABS Internal Report, 2026.

\bibitem{linhares_wolstenholme_2026}
A.~Linhares,
``Deep Vision: A formal proof of Wolstenholme's theorem in Lean~4,''
ARGO LABS Report for Distribution, 2026.

\bibitem{linhares_sketch_2026}
A.~Linhares,
``A sketch of Deep Vision's study of mathematics,''
ARGO LABS Report for Distribution, 2026.

\bibitem{jech_set_theory_2003}
T.~Jech,
\emph{Set Theory: The Third Millennium Edition, Revised and Expanded},
Springer Monographs in Mathematics.
Berlin: Springer-Verlag, 2003.

\bibitem{kunen_set_theory_2011}
K.~Kunen,
\emph{Set Theory: An Introduction to Independence Proofs},
Studies in Logic, vol.~102.
London: College Publications, 2011.

\bibitem{hessenberg_1906}
G.~Hessenberg,
``Grundbegriffe der Mengenlehre,''
\emph{Abhandlungen der Friesschen Schule}, N.S., vol.~1, no.~4,
pp.~485--706, 1906.

\bibitem{tarski_1925}
A.~Tarski,
``Quelques th\'eor\`emes sur les alephs,''
\emph{Fundamenta Mathematicae}, vol.~7, pp.~1--14, 1925.

\bibitem{holz_steffens_weitz_1999}
M.~Holz, K.~Steffens, and E.~Weitz,
\emph{Introduction to Cardinal Arithmetic},
Birkh\"auser Advanced Texts.
Basel: Birkh\"auser Verlag, 1999.

\bibitem{shelah_cardinal_1994}
S.~Shelah,
\emph{Cardinal Arithmetic},
Oxford Logic Guides, vol.~29.
Oxford: Oxford University Press, 1994.

\bibitem{galvin_hajnal_1975}
F.~Galvin and A.~Hajnal,
``Inequalities for cardinal powers,''
\emph{Ann.\ of Math.}, vol.~101, no.~3, pp.~491--498, 1975.

\bibitem{hartogs_1915}
F.~Hartogs,
``\"Uber das Problem der Wohlordnung,''
\emph{Math.\ Annalen}, vol.~76, no.~4, pp.~438--443, 1915.

\bibitem{von_neumann_1923}
J.~von Neumann,
``Zur Einf\"uhrung der transfiniten Zahlen,''
\emph{Acta Litterarum ac Scientiarum Regiae Universitatis
Hungaricae Francisco-Josephinae, Sectio Scientiarum
Mathematicarum}, vol.~1, pp.~199--208, 1923.

\bibitem{drake_1974}
F.~R.~Drake,
\emph{Set Theory: An Introduction to Large Cardinals},
Studies in Logic and the Foundations of Mathematics, vol.~76.
Amsterdam: North-Holland, 1974.

\end{thebibliography}

%======================================================================
% APPENDIX: Full Lean proofs
%======================================================================
\newpage
\appendix

\section{Lean Code: New Proofs}\label{app:lean}

\subsection{Item 13: \texttt{mul\_eq\_left\_iff} via ext1-analog
  (Cardinal.inductionOn + mk\_congr)}

\begin{lstlisting}
import Mathlib.SetTheory.Cardinal.Arithmetic

open Cardinal in
theorem mul_eq_left_iff' {a b : Cardinal} :
    a * b = a <-> max (*$\aleph_0$*) b <= a
      (*$\wedge$*) b (*$\neq$*) 0 (*$\vee$*) b = 1 (*$\vee$*) a = 0 := by
  rw [max_le_iff]
  refine <fun h => ?_, ?_>
  . rcases le_or_gt (*$\aleph_0$*) a with ha | ha
    . have : a (*$\neq$*) 0 := by
        rintro rfl
        exact ha.not_gt aleph0_pos
      left
      rw [and_assoc]
      use ha
      constructor
      . rw [<- not_lt]
        exact fun hb =>
          ne_of_gt (hb.trans_le (le_mul_left this)) h
      . rintro rfl
        apply this
        rw [mul_zero] at h
        exact h.symm
    right
    by_cases h2a : a = 0
    . exact Or.inr h2a
    have hb : b (*$\neq$*) 0 := by
      rintro rfl
      apply h2a
      rw [mul_zero] at h
      exact h.symm
    left
    rw [<- h, mul_lt_aleph0_iff, lt_aleph0,
        lt_aleph0] at ha
    rcases ha with (rfl | rfl | <<n, rfl>, <m, rfl>>)
    . contradiction
    . contradiction
    rw [<- Ne] at h2a
    rw [<- Cardinal.one_le_iff_ne_zero] at h2a hb
    norm_cast at h2a hb h (*$\vdash$*)
    apply le_antisymm _ hb
    rw [<- not_lt]
    apply fun h2b => ne_of_gt _ h
    conv_rhs => left; rw [<- mul_one n]
    rw [Nat.mul_lt_mul_left]
    . exact id
    apply Nat.lt_of_succ_le h2a
  . rintro (<<ha, hab>, hb> | rfl | rfl)
    . rw [mul_eq_max_of_aleph0_le_left ha hb,
          max_eq_left hab]
    -- ext1 analog: reduce cardinal equality to
    -- type equivalence via inductionOn + mk_congr
    . induction a using Cardinal.inductionOn with
      | mk (*$\alpha$*) =>
        rw [show (1 : Cardinal) =
            Cardinal.mk PUnit from
              (Cardinal.mk_eq_one PUnit).symm]
        rw [Cardinal.mul_def]
        exact Cardinal.mk_congr
          (Equiv.prodPUnit (*$\alpha$*))
    . induction b using Cardinal.inductionOn with
      | mk (*$\beta$*) =>
        rw [show (0 : Cardinal) =
            Cardinal.mk PEmpty from
              (Cardinal.mk_eq_zero PEmpty).symm]
        rw [Cardinal.mul_def]
        exact Cardinal.mk_congr
          (Equiv.pemptyProd (*$\beta$*))
\end{lstlisting}

\newpage
\subsection{Item 15: \texttt{mul\_eq\_left\_iff} via
  \texttt{any\_goals lia}}

\begin{lstlisting}
import Mathlib.SetTheory.Cardinal.Arithmetic

open Cardinal

theorem mul_eq_left_iff_reproof {a b : Cardinal} :
    a * b = a <-> max (*$\aleph_0$*) b <= a
      (*$\wedge$*) b (*$\neq$*) 0 (*$\vee$*) b = 1 (*$\vee$*) a = 0 := by
  rw [max_le_iff]
  refine <fun h => ?_, ?_>
  -- [forward direction: identical to Mathlib]
  . rcases le_or_gt (*$\aleph_0$*) a with ha | ha
    . have : a (*$\neq$*) 0 := by
        rintro rfl; exact ha.not_gt aleph0_pos
      left; rw [and_assoc]; use ha
      constructor
      . rw [<- not_lt]
        exact fun hb =>
          ne_of_gt (hb.trans_le (le_mul_left this)) h
      . rintro rfl; apply this
        rw [mul_zero] at h; exact h.symm
    right
    by_cases h2a : a = 0
    . exact Or.inr h2a
    have hb : b (*$\neq$*) 0 := by
      rintro rfl; apply h2a
      rw [mul_zero] at h; exact h.symm
    left
    rw [<- h, mul_lt_aleph0_iff, lt_aleph0,
        lt_aleph0] at ha
    rcases ha with (rfl | rfl | <<n, rfl>, <m, rfl>>)
    . contradiction
    . contradiction
    rw [<- Ne] at h2a
    rw [<- Cardinal.one_le_iff_ne_zero] at h2a hb
    norm_cast at h2a hb h (*$\vdash$*)
    apply le_antisymm _ hb
    rw [<- not_lt]
    apply fun h2b => ne_of_gt _ h
    conv_rhs => left; rw [<- mul_one n]
    rw [Nat.mul_lt_mul_left]
    . exact id
    apply Nat.lt_of_succ_le h2a
  . rintro (<<ha, hab>, hb> | rfl | rfl)
    . rw [mul_eq_max_of_aleph0_le_left ha hb,
          max_eq_left hab]
    -- TRANSFERRED: any_goals lia (from Probability)
    -- Original: all_goals simp
    any_goals lia
\end{lstlisting}

\newpage
\subsection{Item 17: \texttt{mul\_eq\_left\_iff} via top-level
  \texttt{by\_cases ha : $\aleph_0 \leq$ a}}

\begin{lstlisting}
import Mathlib.SetTheory.Cardinal.Arithmetic

open Cardinal

theorem mul_eq_left_iff_reproof {a b : Cardinal} :
    a * b = a <-> max (*$\aleph_0$*) b <= a
      (*$\wedge$*) b (*$\neq$*) 0 (*$\vee$*) b = 1 (*$\vee$*) a = 0 := by
  -- TOP-LEVEL by_cases: structural niceness split
  by_cases ha : (*$\aleph_0$*) <= a
  . -- INFINITE CASE: multiplication = max
    rw [show max (*$\aleph_0$*) b <= a <->
        (*$\aleph_0$*) <= a (*$\wedge$*) b <= a from max_le_iff]
    constructor
    . intro h
      have ha0 : a (*$\neq$*) 0 :=
        ne_of_gt (aleph0_pos.trans_le ha)
      have hb0 : b (*$\neq$*) 0 := by
        rintro rfl; rw [mul_zero] at h
        exact ha0 h.symm
      left
      refine <<ha, ?_>, hb0>
      have h1 :=
        mul_eq_max_of_aleph0_le_left ha hb0
      rw [h] at h1
      exact (max_le_iff.mp (le_of_eq h1.symm)).2
    . rintro (<<_, hab>, hb> | rfl | rfl)
      . exact mul_eq_left ha hab hb
      . simp
      . simp
  . -- FINITE CASE: reduces to Nat arithmetic
    push Not at ha
    constructor
    . intro h
      right
      by_cases h2a : a = 0
      . exact Or.inr h2a
      . left
        have hb : b < (*$\aleph_0$*) := by
          by_contra hb'
          push Not at hb'
          have := mul_eq_max_of_aleph0_le_right
            h2a hb'
          rw [h, max_eq_right
            (ha.trans_le hb').le] at this
          exact ne_of_lt
            (ha.trans_le hb') this
        rw [lt_aleph0] at ha hb
        obtain <n, rfl> := ha
        obtain <m, rfl> := hb
        have hn : n (*$\neq$*) 0 := by
          exact_mod_cast h2a
        norm_cast at h (*$\vdash$*)
        exact (Nat.mul_eq_left hn).mp h
    . rintro (<<h1, _>, _> | rfl | rfl)
      . exact absurd (le_of_max_le_left h1)
          (not_le.mpr ha)
      . simp
      . simp
\end{lstlisting}

\newpage
\subsection{Item 19: \texttt{omega0\_lt\_preOmega\_iff} via
  push-normalization (3 variants)}

\begin{lstlisting}
import Mathlib.SetTheory.Cardinal.Aleph

open Ordinal

-- Variant 1: constructor + rwa
theorem omega0_lt_preOmega_iff_v1
    {x : Ordinal} :
    (*$\omega$*) < preOmega x <-> (*$\omega$*) < x := by
  constructor
  . intro h
    rwa [<- preOmega_omega0,
         preOmega_lt_preOmega] at h
  . intro h
    rwa [<- preOmega_omega0,
         preOmega_lt_preOmega]

-- Variant 2: targeted rewrite
theorem omega0_lt_preOmega_iff_v2
    {x : Ordinal} :
    (*$\omega$*) < preOmega x <-> (*$\omega$*) < x := by
  nth_rw 1 [show ((*$\omega$*) : Ordinal) = preOmega (*$\omega$*)
    from preOmega_omega0.symm]
  exact preOmega_lt_preOmega

-- Variant 3: calc chain
theorem omega0_lt_preOmega_iff_v3
    {x : Ordinal} :
    (*$\omega$*) < preOmega x <-> (*$\omega$*) < x := by
  calc (*$\omega$*) < preOmega x
      <-> preOmega (*$\omega$*) < preOmega x :=
        by rw [preOmega_omega0]
    _ <-> (*$\omega$*) < x := preOmega_lt_preOmega
\end{lstlisting}

\end{document}
